\documentclass[10pt,a4paper]{amsart}
\usepackage[margin=19mm]{geometry}
\usepackage{amsmath,amssymb,mathtools}
\usepackage[scr=rsfs,cal=euler]{mathalfa}
\usepackage{xfrac}
\usepackage[unicode,hidelinks,bookmarks=false]{hyperref}
\newcommand{\ie}{i.e.\ }
\newcommand{\cf}{cf.\ }
\newcommand{\resp}{respectively\ }
\newcommand{\Subsection}{\subsection}
\providecommand{\nobreakdash}{\nobreak\hbox{-}}
\setlength{\emergencystretch}{2em}
\multlinegap0pt
\usepackage{bm}
\usepackage{bbm}
\usepackage{dsfont}
\usepackage{xspace}
\usepackage{cancel}
\usepackage{mathtools}
\usepackage{cleveref}
\usepackage{amsthm}
\makeatletter
\@ifundefined{lem}{\newtheorem{lem}{Lemma}}{}
\makeatother
\makeatletter
\newcommand\ceil[1]{\lceil#1\rceil}
\expandafter\ifx\csname clmproof\endcsname\relax
\newenvironment{clmproof}[1]{\par\noindent\underline{Proof.}\space#1}{\leavevmode\unskip\penalty9999\hbox{}\nobreak\hfill\quad\hbox{$\diamondsuit$}\smallskip}
\fi
\newcommand\floor[1]{\lfloor#1\rfloor}
\renewcommand\ge\geqslant
\renewcommand\le\leqslant
\newcommand\vph\varphi
\expandafter\ifx\csname subclmproof\endcsname\relax
\newenvironment{subclmproof}[1]{\par\noindent\emph{Proof.}\space#1}{\leavevmode\unskip\penalty9999\hbox{}\nobreak\hfill\quad\hbox{$\diamondsuit$}\smallskip}
\fi
\makeatother
\title[Equitably Coloring Planar and Outerplanar Graphs]{Equitably Coloring Planar and Outerplanar Graphs}
\author{Daniel W. Cranston, Reem Mahmoud}
\date{Source: arXiv:2509.16123v2 (CC BY 4.0)}
\begin{document}
\maketitle

\noindent\emph{Source and licence.} Equitably Coloring Planar and Outerplanar Graphs. Version arXiv:2509.16123v2 (CC BY 4.0), distributed under the Creative Commons Attribution 4.0 International licence, as recorded in the arXiv licence metadata for this exact version and shown on the abstract page https://arxiv.org/abs/2509.16123v2. Full rights notice: licenses/arxiv-2509.16123-CC-BY-4.0.txt.\par\smallskip

\noindent\textit{Editorial scope.} Counted unit: the Lem whose source label is \texttt{reduce-to-smallest-s-lem} in the article (source0.tex lines 1446--1450, source label \texttt{reduce-to-smallest-s-lem}), delivered together with the article's own complete original proof of it (source0.tex lines 1452--1621, one \texttt{proof} environment of 170 lines). No mathematics is written, repaired or re-derived: statement and proof are the article's own text line for line. Required context retained uncounted: lem1 (lines 464--472); reduc-to-small-lem (lines 464--472). Every definition, display and label of those lines is retained unchanged. Omitted from this package and not redistributed: 1--463, 473--1445, 1451--1451, 1622--1625. Nothing inside a retained excerpt is omitted or altered: every retained body is delivered exactly as the article's lines are written, so no omission receipt of any kind accompanies this package. Citations: the retained excerpts carry no citation command at all, so no bibliography entry of the article is transcribed and no reference is redistributed. Reference adaptation: none was needed and none was made; every reference inside the delivered unit resolves inside it, so no unresolved reference or citation is produced. Printed slips kept literal: no printed word, symbol, hypothesis or number of the counted statement or of its proof was altered. Layout and class: the article's own class is \texttt{article} with packages including amsmath, amsthm, amssymb, xifthen, fullpage, bm, hyperref, cleveref, verbatim, ulem, algorithmic, algorithm, caption, tikz; the wrapper is the lane's pinned \texttt{amsart} editorial wrapper at 10pt on A4, which restates the theorem-like environments the retained text needs and adds the article's own macro definitions that the retained text needs (\texttt{\textbackslash ceil}, \texttt{\textbackslash floor}, \texttt{\textbackslash ge}, \texttt{\textbackslash le}), with extra wrapper packages (bm, bbm, dsfont, xspace, cancel, mathtools, cleveref). The delivered numbering is the unit's own. Assets: no retained span contains a figure environment, a graphics-inclusion command, a PDF-page inclusion, a table or a TikZ picture; no image file is copied into this package, the package declares no asset dependency and no image file is redistributed with it. Licence: version arXiv:2509.16123v2 of this article is distributed under the Creative Commons Attribution 4.0 International licence, as recorded in the arXiv licence metadata for this exact version and shown on the abstract page https://arxiv.org/abs/2509.16123v2. A full rights notice is delivered at licenses/arxiv-2509.16123-CC-BY-4.0.txt. Characters: every retained excerpt is pure ASCII, so the delivered engine, XeLaTeX with its default Latin Modern text font, reports no missing character.
\begin{lem}
\label{reduce-to-small-s-lem}
\label{lem1}
	Fix an integer $s$ with $s\ge 6$. If $G'$ has an equitable $(s-1)$-coloring whenever $G'$ is outerplanar and 
	$\min_{w\in V(G')}\alpha_{w}(G')\ge \floor{|G'|/(s-1)}$,
	then $G$ has an equitable $s$-coloring whenever $G$ is outerplanar and $\min_{x\in V(G)} \alpha_x(G)\ge \floor{|G|/s}$.
	Thus, for each integer $s_0$ with $s_0\ge 5$, to prove the Main Theorem whenever $s\ge s_0$, it suffices to prove it when $s=s_0$.
\label{reduc-to-small-lem}
\end{lem}

\begin{lem}
\label{reduce-to-smallest-s-lem}
	If every outerplanar graph $H$ with $\min_{w\in V(H)}\alpha_w\ge \floor{|H|/4}$ has an equitable $4$-coloring, 
	then also every outerplanar graph $G$ with $\min_{x\in V(G)}\alpha_x\ge \floor{|G|/5}$ has an equitable $5$-coloring. 
\end{lem}

\begin{proof} 
Let $G$ be an outerplanar graph. 
	By the calculation at the end of \Cref{lem1}, we know that if $(n'+1)/(\alpha'_x+1)>4$, then $d_{G'}(x) > n/5+3.5$. 
	We strengthen this inequality slightly, as follows.  If $d_{G'}(w)\le 3$, then $\alpha'_w\ge 1 + (n'-d_{G'}(w)-1)/4 \ge n'/4$, as desired.
	If $d(w') \ge 4$, then we can strengthen to
	$$
	\left\lceil\frac{n'+1}{\alpha'_w+1}\right\rceil 
	\le \left\lceil\frac{n-\floor{n/5}+1}{\ceil{(n-\floor{n/5}-d(w)-1)/3}+2}\right\rceil
	\le \left\lceil\frac{n-{n/5}+1}{\ceil{(n-{n/5}-d(w)-1)/3}+2}\right\rceil.
	$$
	Note the ``$+1$'' in the final numerator rather than the ``$+2$'' (which holds for all values of $d(w)$) in the proof of \Cref{reduc-to-small-lem}.
	This final inequality holds because $(a+b+\epsilon)/((a+\epsilon)/3) \le (a+b)/(a/3)$ whenever $a,b,\epsilon > 0$.
	With this motivation, let $S:=\{x\in V(G)~|~d_G(x) \ge \ceil{n/5} + 4\}$.  Now we consider $3$ cases based on $|S|$.

\textbf{Case 1: $\bm{|S|\le 1}$.}
Let $w$ be a vertex of maximum degree, and let $I_w$ be an independent set containing $w$ of size $\floor{n/5}$.
Let $G':= G-I_w$.  Clearly $d_{G'}(x)\le n/5+4$ for all $x\in V(G)$.  So $\max_{x\in V(G')}(n'+1)/(\alpha'_x+1)\le 4$, as desired.  

\textbf{Case 2: $\bm{|S|=2}$.}
We denote $S$ by $\{w_1,w_2\}$.  We assume that no independent set $I$ of size $\floor{n/5}$ contains both $w_1$ and $w_2$;
if it does, then we let $G':=G-I$, and we are done.
For each $i\in [2]$, let $I_i$ be an arbitrary independent set containing $w_i$ with size $\floor{n/5}$, and
let $b_i:=|N(w_i)\cap I_{3-i}|$.
If $d(w_i)-b_i\le \ceil{n/5}+3$, for either $i\in [2]$, then we let $G':=G-I_{3-i}$.
Now $w_{3-i}\notin V(G')$ and $d_{G'}(w_i) = d_G(w_i)-|N(w_i)\cap I_{3-i}| = d(w_i) - b_i \le \ceil{n/5}+3$.
Thus, $\max_{x\in V(G')}(n'+1)/(\alpha'_x+1)\le 4$, as desired.  

So assume instead that $d(w_i)-b_i \ge \ceil{n/5}+4$ for each $i\in [2]$.
Let $S_i:=N(w_i)\setminus I_{3-i}$ for each $i\in [2]$.  Since $G$ has no $K_{2,3}$-minor, at most $2$ neighbors of $w_1$ appear on $w_1,w_2$-paths.
Form $S_1'$ from $S_1$ by deleting all such neighbors.  Define $S_2'$ analogously and note, for each $i\in[2]$, that $|S_i'|\ge \ceil{n/5}+2$.
Since $S_1'\subseteq N(w_1)$, we know $G[S_1']$ is $2$-colorable so contains an independent set $I_1'$ of size $|S_1'|/2\ge \ceil{n/10}+1$.
Similarly, $S_2'$ contains an independent set $I_2'$ of size $\ceil{n/10}+1$.  

	By construction $I_1'\cup I_2'$ is also independent; 
	pick $I\subseteq I_1\cup I_2$ with $|I| = \floor{(n+1)/5}$ and let $G':= G-I$.  (Note here $|I|$.  If $\floor{(n+1)/5}=\floor{n/5}$, then we
	know from the calculation above, slightly strengthening that at the end of \Cref{reduc-to-small-lem}, 
	that each other vertex of $G'$ is in an independent set that is sufficiently large.  But if $\floor{(n+1)/5} = \floor{n/5}+1$, then we repeat
	that calculation with $(n+1)/5$ in place of $n/5$.  This time we get $d(w) > n/5+21/5$, which is still good enough.)
	For each $i\in 2$, we know that $\alpha'_{w_i}\ge |I_i| = \floor{n/5} = \floor{n'/4}$.  
Thus, $\max_{x\in V(G')}(n'+1)/(\alpha'_x+1)\le 4$, so we are done.


\textbf{Case 3: $\bm{|S|\ge 3}$.}
We denote the elements of $S$ by $w_i$.  From among all $w_i$ we choose $w_1,w_2,w_3$ (possibly with renaming) to minimize the order of the 
smallest connected subgraph $J$ containing $w_1,w_2,w_3$; we will explain the motivation behind this choice near the end of Cases~3.3 and 3.4.
%
Let $S':=\{w_1,w_2,w_3\}$.  At each mention below of indices $i,j$ or of indices $i,j,k$, we assume that all indices are distinct.
We typically want our statements to hold for all such choices of $i,j$ or of $i,j,k$.
Our goal is to find sets $S_i\subseteq N(w_i)$, for all $i\in [3]$, such that
(1) $|S_i|\ge \lceil n/5\rceil$, (2) $S_i \cap N(w_j) = \emptyset$, and 
(3) $E(S_i,S_j)=\emptyset$.  Note that (2) implies $S_i\cap S_j = \emptyset$.
The value of such sets is that when constructing an independent set of size $\lceil n/5\rceil$ containing $w_i$, we can take $w_i$ along with 
arbitrary independent sets from $G[S_j]$ and $G[S_k]$.  Since $\alpha(G[S_j])\ge |S_j|/2 \ge \lceil n/10\rceil$ and also $\alpha(G[S_k])\ge \lceil 
n/10\rceil$, we will nearly be done.  Finally, we will need to construct big independent sets containing $w_h$ for all $h\ge 4$.

To provide more details, we consider cases based on $G[S]$.
In each case, we start by letting $S_i:=N(w_i)\setminus (N[w_j]\cup N[w_k])$ for each $i\in [3]$.
To ensure condition (3), we may also need to delete a few more vertices from certain sets $S_i$.
Throughout whenever we speak of a $w_i,w_j$-path $P$, we mean that $P$ has length at most $3$ and that $P$ is minimal under inclusion; in particular,
	$|N(w_i)\cap V(P)|=1$ and $|N(w_j)\cap V(P)|=1$.

\textbf{Case 3.1: $\bm{G[S']=K_3}$.}
We focus on ensuring that $E(S_i,S_j)=\emptyset$ in the case $i=1$ and $j=2$, but the same argument works for each choice of distinct $i,j\in[3]$.
Recall that $G$ has no $K_{2,3}$-minor.  Since $w_3\in N(w_1)\cap N(w_2)$, the maximum number of $w_1,w_2$-paths in $G-w_3$ is at most $1$.
(Here the $w_1,w_2$-paths need not be disjoint, just distinct.  If they intersect, we contract this intersection onto its neighbor in 
$\{w_1,w_2\}$.  And we still find a $K_{2,3}$-minor.)  So if such a $w_1,w_2$-path exists, then we simply remove one of its interior
vertices from $S_1$.  We also do the same for $w_1,w_3$-paths.  Thus, in total we remove at most $2$ vertices from our original choice of $S_1$.
So, in the end we have $|S_1|\ge d(w_1)-4 \ge \lceil n/5\rceil$, as desired.  Similarly, we get $|S_2|\ge \lceil n/5\rceil$ and 
$|S_3|\ge \lceil n/5\rceil$.

Since $G$ is outerplanar, $G[S_i]$ is a disjoint union of paths.  We can $2$-color these so that each color class has size at least 
$\lfloor \lceil n/5\rceil/2\rfloor$.  Thus, to construct an independent set $I$ of size $\lfloor n/5\rfloor$ containing $w_1$, we take $w_1$
along with the smaller color class in the $2$-coloring of $G[S_2]$ and the smaller class in the $2$-coloring of $G[S_3]$, 
possibly discarding a vertex to get to size precisely $\lfloor n/5\rfloor$.  Let $G':=G-I$ and $n':=|V(G')|$.  To get an independent set $I'$
in $G'$ containing $w_2$ and having size $\lfloor n'/4\rfloor\le \lceil n/5\rceil$, we take $w_2$ along with the bigger color class in the 
$2$-coloring of $G[S_1]$ and the remaining color class in the $2$-coloring of $G[S_3]$.  
In fact, this argument shows that $\alpha'_{w_2} \ge 2\lceil(\lceil n/5\rceil/2)\rceil + 1 \ge \lceil n/5\rceil+1$.
The same argument works to get an independent set in $G'$ containing $w_3$.

Finally, we consider $w_h$ with $h\ge 4$.  We start with the union of the bigger color class in each $S_i$.  Vertex $w_h$ could have neighbors in
this union, but at most $2$; otherwise we contract $S'$ to a single vertex and get $K_{2,3}$.  So we add $w_h$ to this union and remove its at 
most $2$ neighbors.  The resulting independent set has size at least $3\ceil{n/10}+1-2 \ge \ceil{n/5} \ge \floor{n'/4}$.

\textbf{Case 3.2: $\bm{G[S']=P_3}$.}
By symmetry, we assume that $w_2,w_3\in N(w_1)$.
We can no longer guarantee the sets $S_1,S_2,S_3$ as above, with $|S_i|\ge n/5$ for all $i\in [3]$.
In particular, $w_1$ could have $2$ common neighbors with one or both of $w_2$ and $w_3$.
First suppose that $w_1$ has at most one common neighbor with each of $w_2$ and $w_3$.  Note that $|S_1|\ge d(w_1)-4$.  
For each $w_1,w_2$-path $P$, remove from $S_2$ the neighbor of $w_2$ on $P$ (but remove nothing from $S_1$).
And for each $w_1,w_3$-path $P$, remove from $S_3$ the neighbor of $w_3$ on $P$ (but remove nothing from $S_1$).
As above, it is easy to check that $|S_i|\ge \ceil{n/5}$ for all $i$, so we can finish analogously.

Now assume instead, by symmetry, that $w_1$ has $2$ common neighbors with $w_2$.  In this case, we will ensure that
$|S_1|\ge \ceil{n/5}-2$, but that $|S_2|\ge \ceil{n/5}+1$ and $|S_3|\ge \ceil{n/5}+1$.  Because $w_2$ has $2$ common neighbors with $w_1$, there
do not exist any $w_2,w_3$-paths in $G-N[w_1]$, or we get a $K_{2,3}$-minor.  
	So, $|S_2| = d(w_2) - |N(w_2)\cap N[w_1]| \ge (\ceil{n/5} + 4) - 3 = \ceil{n/5}+1$.
	For each $w_1,w_3$-path $P$, we remove from $S_3$ the neighbor of $w_3$ on $P$.  So similarly, we have $|S_3| \ge d(v) - 3 \ge \ceil{n/5} + 1$.
	And recall that $|S_1| = d(w_1) - |N(w_1)\cap N[w_2]| - |N(w_1)\cap N[w_3]| \ge \ceil{n/5}+4 - 2(3) = \ceil{n/5}-2$.

Now for our independent set $I$ containing $w_1$, we take the smaller color class in each of $G[S_2]$ and $G[S_3]$, along with $w_1$.
This set has size at least $2\lfloor (\lceil n/5\rceil +1)/2\rfloor + 1\ge \lceil n/5\rceil$.  As usual, we let $G':=G-I$.
For an independent set $I'$ in $G'$ of size $\floor{n'/4} = \floor{(n+1)/5}$ containing $w_2$, we take 
$w_2$ along with the bigger color class in each of $G[S_1]$ and $G[S_3]$;
this set has size at least $\lceil(\lceil n/5\rceil + 1)/2\rceil + \lceil(\lceil n/5\rceil - 2)/2\rceil + 1 = 
\lceil(\lceil n/5\rceil + 1)/2\rceil + \lceil\lceil n/5\rceil/2\rceil = \lceil n/5\rceil + 1$.
The same argument works for $w_3$.
Finally, we consider $w_h$ with $h\ge 4$. Similar to above, we start with the bigger color class in each $G[S_i]$.  This set has at most $2$ neighbors
of $w_h$, so we remove them and add $w_h$.  
The resulting set has size at least $\ceil{|S_1|/2}+\ceil{|S_2|/2}+\ceil{|S_3|/2}-2+1\ge (\ceil{n/5}+1)+(\ceil{n/10}-1)-2+1 = \ceil{n/5}+\ceil{n/10}-1 \ge \ceil{n/5}\ge \floor{n'/4}$, as desired.

\textbf{Case 3.3: $\bm{G[S']=\overline{P_3}}$.}
By symmetry, assume that $w_2w_3\in E(G)$ and $w_1w_2,w_1w_3\notin E(G)$.  Note that at most $2$ neighbors of $w_1$ lie on paths from $w_1$ to 
$\{w_2,w_3\}$; otherwise, $G$ has a $K_{2,3}$-minor.  Thus, by removing all such neighbors from $S_1$, we can ensure that $|S_1|\ge d(w_1)-2\ge \ceil{n/5}+2$.
Likewise, since $w_1\notin N(w_2)\cup N(w_3)$, even after removing all neighbors on paths to the other vertices of $S'$, we can ensure that 
$|S_2|\ge \ceil{n/5}-1$ and $|S_3|\ge \ceil{n/5}-1$.  
But actually, we cannot have both of these bounds hold with equality.  
If each of $w_2$ and $w_3$ has both $2$ neighbors on paths to $w_1$ and $2$ neighbors on paths to the other of $w_2,w_3$,
then $G$ contains a $K_{2,3}$-minor.  Thus, we assume by symmetry that $|S_2|\ge \ceil{n/5}$ and $|S_3|\ge \ceil{n/5}-1$.
We construct an independent set $I$ consisting of the smaller color class in each of $G[S_2]$ and $G[S_3]$, along with $w_1$.
This set has size at least $\floor{(\ceil{n/5} -1)/2} + \floor{\ceil{n/5}/2} + 1 = (\ceil{n/5}-1)+1 = \ceil{n/5}$.  
Let $G':= G-I$.
In $G'$, we construct an independent set $I'$ consisting of the bigger color class in $G[S_1]$, the bigger color class in $G[S_3]$, 
along with $w_2$.  This set has size at least 
$\ceil{(\ceil{n/5} + 2)/2} + \ceil{(\ceil{n/5}-1)/2} + 1=
\ceil{\ceil{n/5}/2} + \ceil{(\ceil{n/5} +1)/2} + 1=
\ceil{n/5}+ 2 \ge \floor{n'/4}$. 
And the independent set for $w_3$ is similar.

Finally, we consider $w_h$ with $h\ge 4$.  Now we start with the independent set above in $G'$ containing $w_2$, remove $w_2$, add $w_h$, add the bigger color class in $G[S_2]$, and remove all neighbors of $w_h$.  Before removing neighbors of $w_h$, this set has size at least $(\ceil{n/5}+2)+\ceil{\ceil{n/5}/2}$. 
Clearly, $\ceil{\ceil{n/5}/2}\ge 1$.
So it suffices to show that $w_h$ has at most $3$ neighbors in this set; suppose not.  We now contract the smallest connected subgraph $J$ containing
all of $S'$ to a single vertex, and get a copy of $K_{2,3}$.  Note that $w_h\notin N(w_2)\cup N(w_3)$, by our choice of $w_1,w_2,w_3$ to minimize the 
order of $J$.  First suppose that $w_h$ has a neighbor in $S_2\cup S_3$.  This implies that $w_1$ must have a neighbor in $N(w_2)\cup N(w_3)$.
So $|V(J)|=4$.  Thus, at most one neighbor of $w_h$ in $S_1\cup S_2\cup S_3$ is contracted away; so $3$ remaining neighbors of $w_h$ give the part of 
size $3$ in the copy of $K_{2,3}$.  So assume instead that $w_h$ has no neighbor in $S_2\cup S_3$.  Thus, all neighbors of $w_h$ in $S_1\cup S_2\cup S_3$
lie in $S_1$.  But this must be at most $2$ such neighbors, since $G$ has no $K_{2,3}$-minor, a contradiction.

\textbf{Case 3.4: $\bm{G[S']=\overline{K_3}}$.}
This case is similar to the previous ones, but we need to prove stronger lower bounds on $|S_i|$ to allow that $w_h$ might have up to $4$ neighbors in
$S_1\cup S_2\cup S_3$.  
If, for every pair $w_i,w_j$ the subgraph $G-N[w_k]$ contains a $w_i,w_j$ path, then each $N(w_i)$ contains at most $2$ vertices on $w_i,\{w_j,w_k\}$-paths.
We remove all of these and get $|S_i|\ge \ceil{n/5}+2$ for all $i$.  So assume instead, by symmetry, that $G$ contains no such $w_1,w_2$-path. 
If $G$ contains no $w_1,w_2$-path of length at most $3$, then even after removing all neighbors of $w_i$ on paths to $\{w_j,w_k\}$, for each choice of $i,j,k$,
we have $|S_1|\ge \ceil{n/5}+2$, $|S_2|\ge \ceil{n/5}+2$, and $|S_3|\ge \ceil{n/5}$.  So assume instead that such a path $P$ exists.  If $P$ has length $2$,
then its interior vertex $x$ lies in $N(w_1)\cap N(w_2)\cap N(w_3)$.  Furthermore, every $w_1,w_2$-path of length at most $3$ contains $x$.  Now removing
from $S_i$ all neighbors of $w_i$ on $w_i,w_3$-paths, for each $i\in [2]$, we still have 
$|S_1|\ge \ceil{n/5}+2$, $|S_2|\ge \ceil{n/5}+2$, and $|S_3|\ge \ceil{n/5}$.
So we assume below that $N(w_1)\cap N(w_2)\cap N(w_3) = \emptyset$.


So assume instead that (no such path $w_1,w_2$-path of length $2$ exists and) $P$ has length $3$.
Denote $P$ by $w_1x_1x_2w_2$.  By symmetry, assume that $x_2\in N(w_3)$.  So $w_1$ has at most $2$ neighbors on $w_1,\{w_2,w_3\}$-paths.
If $G$ also contains a $w_1,w_3$-path disjoint from $N(w_2)$, then also $w_2$ has at most $2$ neighbors on paths to $\{w_1,w_3\}$,
so $|S_1|\ge \ceil{n/5}+2$, $|S_2|\ge \ceil{n/5}+2$, and $|S_3|\ge \ceil{n/5}$.  So assume that $G$ contains no such path.
We will show that every $w_3,w_1$-path contains $x_2$.  Suppose not.  Consider a $w_3,w_1$-path $P$ of length $3$, and let $x_3$
be an internal vertex of $P$ such that $x_3\ne x_2$ and $x_3\in N(w_2)$.  Now $G$ has a $K_{2,3}$-minor with $\{x_2,x_3\}$ as one part and 
with $\{w_1,w_2,w_3\}$ as the other part.  Thus, in every case we get $|S_1|\ge \ceil{n/5}+2$, $|S_2|\ge \ceil{n/5}+2$, and $|S_3|\ge \ceil{n/5}$.
For our independent set $I$, we take $w_1$, along with the smaller color class in $G[S_2]$ and the smaller color class in $G[S_3]$.
This gives an independent set of size bigger than $\ceil{n/5}$, and we take $I$ to be an arbitrary subset of size $\floor{n/5}$.  
Let $G':=G-I$.  The analysis is straightforward for independent sets in $G'$ of size $\floor{n'/4}\le \ceil{n/5}$
that contain $w_2$ or contain $w_3$.

Finally, we consider $w_h$ with $h\ge 4$.  We start with the independent set containing the bigger color class in each of $G[S_1]$, $G[S_2]$, and $G[S_3]$.
Now we add $w_h$ and remove all neighbors of $w_h$.  The size of this set before removing neighbors of $w_h$ is at least
$2\ceil{(\ceil{n/5}+2)/2}+\ceil{\ceil{n/5}/2}+1\ge \ceil{n/5}+\ceil{n/10}+3\ge \ceil{n/5}+4$.  So it suffices to show that $w_h$ has at most $4$ neighbors
in the set.  If $w_h$ has more than $4$ neighbors in $S_1\cup S_2\cup S_3$, then it has neighbors in each $S_i$ (since it has at most $2$ in each).
This implies that $|V(J)|\le5$; otherwise, we could get a smaller $J$ by replacing some $w_i$ with $w_h$.  Now we contract $J$ to a single vertex.
Doing so contracts away at most $2$ neighbors of $w_h$.  Thus, the remaining at least $3$ neighbors of $w_h$ in $S_1\cup S_2\cup S_3$ give rise to a 
copy of $K_{2,3}$, a contradiction.
\end{proof}

\end{document}
